\centerline{\bf \Huge Домашнее задание}
\centerline{\bf \Huge Степень вхождения}

\begin{flushright}
        \scriptsize{Если ты первый, ты первый\\ Илья Михайлович}
\end{flushright}

\problem{1} 
Натуральные числа $a$ и $b$ таковы, что сумма $\dfrac{a^2}{b} + \dfrac{b^2}{a}$ целая. Докажите, что оба слагаемых целые.

\problem{2} 
Можно ли из чисел $1!,2!,...,99!,100!$   вычеркнуть одно так, чтобы произведение оставшихся оказалось кубом натурального числа?


\problem{3} 
Даны натуральные числа $a$ и $b$ такие, что число $\frac{a+1}{b}+\frac{b+1}{a}$ является целым. Докажите, что наибольший общий делитель чисел $a$ и $b$ не превосходит числа $\sqrt{a+b}$.

\problem{4} 
Существует ли бесконечная арифметическая прогрессия с ненулевой разностью, состоящая только из степеней (выше первой) натуральных чисел?

\problem{5}
Натуральные числа $a, b, c$ таковы, что число $\dfrac{a}{b} + \dfrac{b}{c} +\dfrac{c}{a}$ является целым. Верно ли, что  $abc \ - $ точный куб?
